\section{Related Work}

\paragraph{LoRA adaptation and adapter risk.}
LoRA reduces fine-tuning cost by representing updates through low-rank matrices injected into selected transformer projections \citep{hu2022lora}.
Its modularity makes adapters easy to train, distribute, merge, and remove.
That same modularity also creates an attack surface: the base model may remain unchanged while a small adapter controls downstream behavior.
Recent work on LLM backdoors and sleeper behavior shows that malicious behaviors can persist through standard alignment or fine-tuning procedures \citep{wang2024sleeper,zhang2024backdoor,bairi2024backdoor}.
Our work focuses on the adapter setting and asks whether deployment-time compression can selectively disrupt the malicious behavior.

\paragraph{LLM compression.}
Compression methods such as pruning, quantization, and low-rank approximation are usually evaluated through efficiency and utility \citep{han2015deep,frantar2023gptq,gao2024compression}.
For backdoored adapters, however, utility alone is not enough.
Quantization may preserve the harmful pathway; pruning may suppress the attack only after inducing broad refusal.
We therefore treat compression as a behavioral intervention whose security effect must be measured directly.

\paragraph{Backdoor defenses.}
Backdoor defenses for language models include data filtering, representation editing, safety tuning, and post-hoc detection \citep{qi2023bad,li2023bad,wang2024vaccine}.
Many defenses operate on activations, prompts, or training data.
\sac{} instead acts on the adapter weights themselves.
It is closest in spirit to post-training model editing: the goal is to remove a behavior without retraining the full model.

\paragraph{Mechanistic and causal analysis.}
Causal tracing and mechanistic interpretability aim to localize internal representations responsible for model behavior \citep{meng2022locating,elhage2021mathematical,yang2024mechanistic,cunningham2023sparse}.
Mechanistic analysis motivates targeted interventions, but our focus is the deployment-time compression problem.
Rather than claiming a universal mechanistic circuit, we use behaviorally grounded counterfactual metrics to select compression operations and then test whether the resulting adapter improves the safety-utility frontier.
